\clearpage
\begin{center}
{\Large DETERMINATION}\par\vspace{8mm}
{\small OF THE}\par\vspace{8mm}
{\LARGE\textls[100]{DIFFERENCE IN LATITUDE}}\par\vspace{8mm}
{\small BETWEEN THE}\par\vspace{7mm}
{\large OBSERVATORIES OF GÖTTINGEN AND ALTONA}\par\vspace{7mm}
{\small BY}\par\vspace{7mm}
{\normalsize OBSERVATIONS WITH RAMSDEN'S ZENITH SECTOR}\par\vspace{7mm}
{\small BY}\par\vspace{7mm}
{\large\textls[100]{CARL FRIEDRICH GAUSS,}}\par\smallskip
{\scriptsize
KNIGHT OF THE GUELPHIC ORDER AND THE ORDER OF THE DANNEBROG;\\
ROYAL BRITISH-HANOVERIAN COURT COUNCILLOR;\\
PROFESSOR OF ASTRONOMY AND DIRECTOR OF THE OBSERVATORY IN GÖTTINGEN;\\
MEMBER OF THE ACADEMIES AND SOCIETIES OF BERLIN, COPENHAGEN, EDINBURGH, GÖTTINGEN,
LONDON, MUNICH, NAPLES, PARIS, ST PETERSBURG, AND STOCKHOLM, AND OF THE
AMERICAN, ITALIAN, COURLAND, AND LONDON ASTRONOMICAL SOCIETIES, AMONG OTHERS.
}\par\vspace{15mm}
\rule{.35\textwidth}{.4pt}\par\medskip
{\large\textls[170]{GÖTTINGEN,}}\par\medskip
PUBLISHED BY VANDENHOECK AND RUPRECHT.\par\medskip
1828.
\end{center}
\vfill\noindent IX.\hfill 1\par

\clearpage
\null

\clearpage
\begin{center}
{\LARGE\textls[170]{CONTENTS.}}\par\medskip
\rule{.09\textwidth}{.4pt}
\end{center}
\noindent Introduction\dotfill p.\quad 5\par
\noindent\makebox[8mm][r]{I.}\quad The Stars Observed.\dotfill \quotedblbase\quad 8\par
\noindent\makebox[8mm][r]{II.}\quad The Observations\dotfill \quotedblbase\quad 10\par
\noindent\makebox[8mm][r]{III.}\quad Results\dotfill \quotedblbase\quad 29\par
\noindent\hspace*{12mm}\begin{minipage}{\dimexpr\textwidth-24mm\relax}
Simplest combination of the observations for determining the difference in latitude, Art.\ 1.
Accuracy of the observations, Art.\ 2. Collimation error, Art.\ 3.
Absolutely most advantageous combination of the observations, Art.\ 4---7.
Allowance for irregular errors of graduation, Art.\ 8---11. Positions of the observing sites, Art.\ 12.
Determination of the absolute altitude of the pole at the Göttingen Observatory from observations
of the Pole Star with \textsc{Reichenbach}'s meridian circle, Art.\ 13---17.
Final result of the Hanoverian arc measurement, Art.\ 18---20.
Comparison of the declinations of the zenith stars observed with the determinations
of \textsc{Bradley} and \textsc{Piazzi}, Art.\ 21.
\end{minipage}\par
\noindent\makebox[8mm][r]{IV.}\quad Determination of the Latitude of the Seeberg Observatory\dotfill \quotedblbase\quad 52\par
\noindent Supplement to [Art.\ 20] p.\ 48\dotfill \quotedblbase\quad 56\par
\begin{center}\rule{.23\textwidth}{.4pt}\end{center}
\vfill\begin{flushright}$1^{\ast}$\end{flushright}

\clearpage
\null

\clearpage
\vspace*{12mm}
\begin{center}
{\large DETERMINATION OF THE DIFFERENCE IN LATITUDE}\par\medskip
{\small BETWEEN}\par\medskip
{\large THE OBSERVATORIES OF GÖTTINGEN AND ALTONA\\
BY OBSERVATIONS\\
WITH RAMSDEN'S ZENITH SECTOR.}\par\medskip
\rule{.12\textwidth}{.4pt}\par\bigskip
{\large INTRODUCTION.}
\end{center}

The chain of triangles that I ran through the Kingdom of Hanover along the Göttingen meridian
in the years 1821---1824 connects the observatories of Göttingen and Altona trigonometrically
with the utmost precision. These measurements will be published in detail in due course.
I note here only that the absolute magnitudes are based on the baseline measured in Holstein
with the utmost precision by Professor \textsc{Schumacher}, to which the system of triangles
is connected by the Hamburg-Hohenhorn side. The orientation is based on observations with the
Göttingen transit telescope, since the observatory and the northern meridian mark are themselves
triangulation points. By a remarkable coincidence, the observatories of Göttingen and Altona
lie on the same meridian to within less than the width of a house. Although the absolute
altitudes of the pole have been determined by observations with fixed meridian instruments,
it was nevertheless important to determine the difference in latitude by another method with
one and the same instrument, and I was fortunate

\clearpage
\noindent\makebox[10mm][l]{6}\parbox[t]{\dimexpr\textwidth-10mm\relax}{\centering\scriptsize DETERMINATION OF THE\,\quotedblbase\,DIFFERENCE IN LATITUDE BETWEEN GÖTTINGEN AND ALTONA.}\par
\noindent to be able to use for this purpose the excellent zenith sector of \textsc{Ramsden},
which, as is well known, was used for a similar purpose in the English arc measurement.
The observations I made with it in the spring of 1827 and their results are the principal
subject of this work.

When many stars have to be observed in succession, observations with this instrument can
hardly be made without the help of a practised assistant. Professor \textsc{Schumacher} was
therefore kind enough, with the consent of His Majesty the King of Denmark, to assign Engineer
Lieutenant von \textsc{Nehus}[\textsuperscript{*})] to the observations at both sites. This very
skilful observer continuously attended to the reading of the micrometer screw and the adjustment
of the plumb line, while I myself observed the crossings of the meridian wires and set the wire
perpendicular to the meridian on the stars. Only on the first two nights of observation at
Altona was the former task performed by another assistant. Those observations have therefore
been left out, especially since experience confirmed that different people judged the bisection
of the points by the plumb line differently.

The instrument is sufficiently well known from \textsc{Mudge}'s detailed description.
In Göttingen it could be set up in the observatory itself, under the eastern meridian slit.
In Altona this was impracticable. It was therefore set up in Professor \textsc{Schumacher}'s
garden, where the observatory itself stands, under the same observing tent that
\textsc{Mudge} used in England. The stability of the mounting on driven piles left nothing
to be desired. The levelling of the vertical axis was checked daily, and usually almost nothing
was found to need adjustment. The same applies to the horizontal axis.

To bring the plane of the limb into the meridian, we used the southern meridian mark in
Göttingen. It stands in the meridian of the western slit, but its azimuth at the site of the
sector could be calculated with the greatest precision. No similar means could be used in
Altona. There the limb was first brought, with the help of the known absolute time and by means
of a culminating star, very nearly into the

\vfill\noindent\rule{.24\textwidth}{.4pt}\par
{\footnotesize\noindent [\textsuperscript{*}) Manuscript note:] \textsc{Nehus} died of influenza that had receded, 17 April 1844.\par}

\clearpage
\noindent\mbox{}\hfill INTRODUCTION.\hfill 7\par
\noindent meridian. Observing several stars over the full extent of the limb then readily gave
the small correction still needed for the setting. As already mentioned, the crossings of the
meridian wires by the stars culminating at the sector were also observed each night, and the
right ascensions of the stars were known. This afforded a continual and reliable check of
the correct setting in the meridian, and only once was a slight readjustment necessary.
As a rule the position of the limb was alternated from night to night, east or west. It was
only toward the end of the observations, in order to distribute the number of observations
fairly evenly between the two positions, that we sometimes departed from this rule and turned
the sector once or several times in a single night.

The readings of the barometer and of the inside and outside thermometers were recorded at
least three times each night, at the beginning, middle, and end of the observations.
Following \textsc{Mudge}'s example, the difference in temperature at the top and bottom of
the sector was also recorded, since it changes the limb and the radius in unequal proportions.
It hardly needs separate mention that every other precaution required by the design of the
instrument was carefully observed, for example, keeping the vessel of water in which the plumb
bob hangs properly filled, keeping the same threads of the micrometer screw in play as far as
possible, and the like. For each star the plumb line was set afresh on the upper point
(the centre of the graduated arc), independently of the preceding setting. The setting on the
nearest division (or the two nearest divisions) was generally repeated several times, and the
mean was taken of the several readings of the micrometer screw, which usually agreed to within
a few decimal fractions of a second.

\begin{center}\rule{.24\textwidth}{.4pt}\end{center}

\clearpage
\vspace*{18mm}
\begin{center}
{\large I.}\par\medskip
{\Large THE STARS OBSERVED.}
\end{center}

At first I had selected 38 stars suitably placed for observation. Toward the end of the
observations at Göttingen I added five others, fearing that unfavourable weather might delay
the end of the observations at Altona so much that a considerable number of the original stars,
now culminating by daylight, could not be observed often enough. This fear proved justified
only to a small degree, and only one star was observed at Altona on one side alone. I give here
the mean positions of these stars, reduced to the beginning of 1827. The declinations are the
results yielded by the observations with the zenith sector itself. For the right ascensions,
which need not be determined with the utmost precision for the present purpose, the basis
is usually only a single observation with the meridian circle, whose reduction Mr.\
\textsc{von Heiligenstein} kindly calculated. For convenience I designate the stars by
consecutive numbers. Nos.\ 8, 13, 15, and 31 are double stars. For the first the following star
was always set on the wire, and for the next two the midpoint. In No.\ 31 the companion is so
faint that it remained invisible in the sector telescope, even when the field was darkened
as a trial. Professor \textsc{Schumacher}, however, noticed it immediately in the brighter
telescope of the \textsc{Reichenbach} meridian circle, although at the time we did not know
that other astronomers had already recognized this star as double.

\clearpage
\noindent\parbox[t]{\dimexpr\textwidth-10mm\relax}{\centering\scriptsize DETERMINATION OF THE DIFFERENCE IN LATITUDE ETC.\quad I.\quad THE STARS OBSERVED.}\makebox[10mm][r]{9}\par
\begin{center}
\begin{minipage}[t]{.49\textwidth}\centering
{\scriptsize\setlength{\tabcolsep}{1.2pt}\renewcommand{\arraystretch}{1.2}\begin{tabular}{r|l|rrr|rrr|}
\multicolumn{2}{c|}{Designation}&\multicolumn{3}{c|}{R.A. 1827}&\multicolumn{3}{c|}{Decl. 1827}\\\hline
1&24 Canum&$13^{\mathrm h}$&$27^{\mathrm m}$&$22^{\mathrm s}{,}39$&$49^\circ$&$54'$&$11''{,}62$\\
2&83 Ursae&&$34$&$9{,}77$&$55$&$33$&$34{,}98$\\
3&$\eta$ Ursae&&$40$&$42{,}86$&$50$&$10$&$46{,}20$\\
4&86 Ursae&&$47$&$28{,}48$&$54$&$34$&$55{,}55$\\
5&&&$50$&$45{,}55$&$55$&$25$&$59{,}22$\\
6&P. 13. 289&&$55$&$19{,}44$&$46$&$35$&$39{,}14$\\
7&13 Bootis&$14$&$1$&$49{,}02$&$50$&$16$&$43{,}44$\\
8&$\varkappa$ Bootis seq.&&$7$&$16{,}93$&$52$&$36$&$7{,}47$\\
9&P. 14. 56&&$12$&$5{,}39$&$56$&$13$&$36{,}51$\\
10&$\theta$ Bootis&&$19$&$18{,}40$&$52$&$39$&$12{,}05$\\
11&P. 14. 131&&$27$&$50{,}42$&$53$&$39$&$33{,}02$\\
12&P. 14. 164&&$35$&$24{,}18$&$52$&$58$&$55{,}66$\\
13&39 Bootis med.&&$43$&$48{,}37$&$49$&$26$&$9{,}46$\\
14&P. 14. 235&&$50$&$38{,}71$&$50$&$20$&$22{,}49$\\
15&44 Bootis med.&&$58$&$5{,}30$&$48$&$19$&$52{,}47$\\
16&&$15$&$7$&$8{,}12$&$49$&$13$&$46{,}68$\\
17&P. 15. 39&&$10$&$33{,}95$&$51$&$34$&$53{,}72$\\
18&&&$15$&$0{,}55$&$52$&$35$&$5{,}87$\\
19&&&$21$&$50{,}88$&$54$&$37$&$34{,}87$\\
20&&&$30$&$42{,}13$&$54$&$29$&$54{,}30$\\
21&&&$38$&$9{,}07$&$52$&$54$&$37{,}09$\\
22&&&$42$&$14{,}32$&$46$&$16$&$4{,}70$\\
\end{tabular}}\end{minipage}\hfill
\begin{minipage}[t]{.49\textwidth}\centering
{\scriptsize\setlength{\tabcolsep}{1.2pt}\renewcommand{\arraystretch}{1.2}\begin{tabular}{r|l|rrr|rrr}
\multicolumn{2}{c|}{Designation}&\multicolumn{3}{c|}{R.A. 1827}&\multicolumn{3}{c|}{Decl. 1827}\\\hline
23&&$15^{\mathrm h}$&$48^{\mathrm m}$&$16^{\mathrm s}{,}24$&$56^\circ$&$20'$&$27''{,}00$\\
24&&&$54$&$10{,}23$&$50$&$22$&$39{,}41$\\
25&$\theta$ Draconis&&$58$&$39{,}70$&$59$&$1$&$47{,}14$\\
26&&$16$&$4$&$24{,}06$&$50$&$38$&$13{,}12$\\
27&P. 16. 33&&$7$&$24{,}92$&$46$&$20$&$17{,}04$\\
28&P. 16. 56&&$11$&$34{,}31$&$53$&$40$&$13{,}41$\\
29&&&$16$&$2{,}19$&$52$&$27$&$11{,}06$\\
30&&&$20$&$38{,}75$&$55$&$36$&$4{,}73$\\
31&&&$26$&$34{,}90$&$45$&$58$&$5{,}85$\\
32&16 Draconis&&$32$&$6{,}42$&$53$&$15$&$3{,}64$\\
33&&&$37$&$53{,}85$&$50$&$16$&$13{,}67$\\
34&&&$42$&$1{,}51$&$57$&$5$&$37{,}79$\\
35&&&$45$&$2{,}26$&$46$&$56$&$48{,}53$\\
36&P. 16. 253&&$49$&$21{,}51$&$46$&$49$&$22{,}01$\\
37&P. 16. 291&&$56$&$11{,}73$&$56$&$56$&$42{,}35$\\
38&P. 16. 310&$17$&$0$&$15{,}10$&$49$&$2$&$47{,}09$\\
39&P. 17. 20&&$4$&$23{,}12$&$58$&$29$&$49{,}33$\\
40&P. 17. 38&&$8$&$1{,}46$&$56$&$52$&$21{,}74$\\
41&74 Herculis&&$15$&$28{,}30$&$46$&$24$&$52{,}90$\\
42&P. 17. 120&&$20$&$53{,}56$&$57$&$10$&$13{,}11$\\
43&$\beta$ Draconis&&$26$&$31{,}90$&$52$&$25$&$57{,}91$\\
&&&&&&&\\
\end{tabular}}\end{minipage}
\end{center}
\begin{center}\rule{.24\textwidth}{.4pt}\end{center}

\clearpage
\vspace*{18mm}
\begin{center}
{\large II.}\par\medskip
{\Large THE OBSERVATIONS.}
\end{center}

A complete printing of the journal in its original form, which would have more than doubled
the length of this work, seemed unnecessary to me. I therefore give the observations at once
in order by star.

The first column contains the zenith distances as given by the instrument, that is, the mere
reduction of the reading. Northern zenith distances are taken as positive, southern ones as negative.

The second column gives the combined effect of refraction and the unequal expansion of the
instrument owing to the difference in temperature between its upper and lower parts. The extreme
differences that occurred were $+1^\circ{,}2$ Réaum.\ (the upper thermometer reading higher) and
$-0^\circ{,}6$. To obtain a result suitable for judging how well the observations agree with
one another, I took the trouble to compute the amount for each individual observation, taking
advantage of a few minor computational shortcuts that readily presented themselves.

The third column contains the reduction to the mean position at the beginning of the year
for aberration, nutation, and precession. For some stars the proper motion is added as well.
Specifically, the annual proper motion in declination is taken to be
\begin{center}
\begin{tabular}{@{}lr@{\quad}c@{\quad}l@{}}
for&10&.\ .\ .\ .\ .\ .\ .&$-0''{,}42$\\
&25&.\ .\ .\ .\ .\ .\ .&$+0{,}33$\\
&37&.\ .\ .\ .\ .\ .\ .&$+0{,}38$.
\end{tabular}
\end{center}

\clearpage
\noindent\parbox[t]{\dimexpr\textwidth-10mm\relax}{\centering\scriptsize DETERMINATION OF THE DIFFERENCE IN LATITUDE ETC.\quad II.\quad THE OBSERVATIONS.}\makebox[10mm][r]{11}\par
\noindent For the first two stars the proper motion has long been established. In No.\ 37 it
emerges from comparison with \textsc{Piazzi}'s determination in such a way that, assuming the
latter is correct, it cannot be doubted\textsuperscript{*}). \textsc{Baily}'s valuable tables
are the basis of the calculation of aberration, nutation, and precession. From them an
ephemeris at ten-day intervals was calculated for each star with the help of Messrs.\
von \textsc{Nehus} and \textsc{Petersen}, and interpolation was carried out in it with allowance
for the second differences.

Finally, the fourth column contains the sum of the first three, hence the true zenith
distance, still affected only by the collimation error, for the mean position at the
beginning of 1827 as obtained from each individual observation.

\medskip\noindent\rule{.24\textwidth}{.4pt}\par
{\footnotesize
\noindent\textsuperscript{*}) The correctness of \textsc{Piazzi}'s determination of this star,
based on 8 observations, is corroborated by its close agreement with the value in the earlier
edition of his catalogue, from 1803, which was based on 6 observations. The exact magnitude of
the proper motion nevertheless remains somewhat uncertain, since the year corresponding to the
mean of the observations is unknown. It is remarkable to find this considerable proper motion
in a star of the 7th magnitude. No.\ 11 also seems to deserve the attention of astronomers in
this respect.\par}
\begin{center}\rule{.24\textwidth}{.4pt}\end{center}
